\documentclass[11pt]{article}

\usepackage[utf8]{inputenc}
\usepackage[T1]{fontenc}
\usepackage{lmodern}
\usepackage{amsmath,amssymb,amsthm}
\usepackage{geometry}

\geometry{margin=1in}

\title{Short-cycle Schur corrections in graphon theory}
\author{Taeyoung Kim}
\date{\today}

\begin{document}

\maketitle

\subsection{The regular case and the short-cycle Schur corrections}
\label{subsec:regular-short-schur}

The elementary arguments of Section~2 are shorter than the general machinery developed here.
The purpose of this subsection is different: it shows that the cases
$m=3,5,7$ arise as the first three instances of the same regular-to-general
mechanism used in the remainder of the proof.

Suppose first that $U$ is $q$-regular, or equivalently that $W$ is
$p$-regular. Then $g=0$, and hence $\mu=0$, $L_-=L_+=0$, and $S_m=0$.
Moreover,
\[
 x_j=\langle \mathbf 1,T_U^j\mathbf 1\rangle=q^j
 \qquad (j\ge 0).
\]
Consequently, the reduction of Sections~4.1--4.2 gives
\[
 t(C_m,W)-\bigl(p^m-pq^{m-1}\bigr)
 \ge q^{m-1}-x_{m-1}=0.
\]
Thus the regular case is exactly the case in which no excursion from the
constant direction to $\mathbf 1^\perp$ occurs.

For a nonregular graphon, Proposition~4.2 writes the error term $\Phi_m$ as
an integral of the kernels $P_{m,r}$ against the positive measures
$d\mu^{\otimes r}$.  The index $r$ records the number of excursions from the
constant direction into the compressed operator $A$ and back.  For the first
three odd cycles, each of these excursion kernels is already pointwise
nonnegative.

\begin{proposition}[The excursion kernels for $C_3,C_5,C_7$]
\label{prop:short-cycle-excursion-kernels}
Although the standing assumption in this section is $q\le 1/3$, the
identities of Sections~4.1--4.3 are algebraic and the following calculation
is valid throughout the full nontrivial range $0\le q\le 1/2$.
For $m\in\{3,5,7\}$, every
$1\le r\le (m-1)/2$, and every choice of real
$\lambda_1,\ldots,\lambda_r$, one has
\[
 P_{m,r}(q;\lambda_1,\ldots,\lambda_r)\ge 0.
\]
Moreover, $P_{m,1}(q;\lambda)>0$ for every real $\lambda$.  Consequently,
\[
 t(C_m,W)\ge p^m-pq^{m-1}
 \qquad (m=3,5,7),
\]
and the inequality is strict whenever $g\ne 0$, equivalently whenever $W$
is not regular.
\end{proposition}

\begin{proof}
For $m=3$, only $r=1$ occurs, and the definition of $P_{m,r}$ gives
\[
 \begin{aligned}
 P_{3,1}(q;\lambda)
 &=3\bigl[h_1(p,-\lambda)+h_1(q,\lambda)\bigr]-h_0(q,q,\lambda)\\
 &=3\bigl[(p-\lambda)+(q+\lambda)\bigr]-1=2.
 \end{aligned}
\]
Hence
\[
 \Phi_3=2\mu(\mathbb R)=2\lVert g\rVert_2^2.
\]

For $m=5$, the two possible excursion numbers give
\[
 P_{5,1}(q;\lambda)
 =4\lambda^2+(8q-5)\lambda+12q^2-15q+5
\]
and
\[
 P_{5,2}(q;\lambda_1,\lambda_2)=4.
\]
The first polynomial has the sum-of-squares decomposition
\[
 P_{5,1}(q;\lambda)
 =4\left(\lambda+q-\frac58\right)^2
  +8\left(q-\frac58\right)^2+\frac5{16},
\]
so it is strictly positive for every $q$ and every real $\lambda$.  Therefore
\[
 \Phi_5
 =\int P_{5,1}(q;\lambda)\,d\mu(\lambda)
  +4\mu(\mathbb R)^2
 \ge 0.
\]

For $m=7$, there are three possible excursion numbers.  For $r=1$ one
obtains
\[
 \begin{aligned}
 P_{7,1}(q;\lambda)={}&
 6\lambda^4+(12q-7)\lambda^3
 +(18q^2-21q+7)\lambda^2\\
 &+(24q^3-42q^2+28q-7)\lambda\\
 &+30q^4-70q^3+70q^2-35q+7.
 \end{aligned}
\]
This is exactly the quartic $P_q(\lambda)$ treated in Section~2.4.  The
sum-of-squares decomposition established there shows that
\[
 P_{7,1}(q;\lambda)>0
 \qquad
 \left(0\le q\le\frac12,\ \lambda\in\mathbb R\right).
\]
For $r=2$, direct expansion gives
\[
 \begin{aligned}
 P_{7,2}(q;\lambda_1,\lambda_2)={}&
 6(\lambda_1^2+\lambda_1\lambda_2+\lambda_2^2)
 +\left(18q-\frac{21}{2}\right)(\lambda_1+\lambda_2)\\
 &+36q^2-42q+14,
 \end{aligned}
\]
and completing squares yields
\[
 \begin{aligned}
 P_{7,2}(q;\lambda_1,\lambda_2)={}&
 \frac32(\lambda_1-\lambda_2)^2
 +\frac92\left(\lambda_1+\lambda_2+2q-\frac76\right)^2\\
 &+18\left(q-\frac7{12}\right)^2+\frac74.
 \end{aligned}
\]
Finally,
\[
 P_{7,3}(q;\lambda_1,\lambda_2,\lambda_3)=6.
\]
It follows that
\[
 \begin{aligned}
 \Phi_7={}&
 \int P_{7,1}(q;\lambda)\,d\mu(\lambda)\\
 &+\iint P_{7,2}(q;\lambda_1,\lambda_2)
 \,d\mu(\lambda_1)d\mu(\lambda_2)
 +6\mu(\mathbb R)^3
 \ge 0.
 \end{aligned}
\]

In all three cases, if $g\ne0$, then $\mu(\mathbb R)=\lVert g\rVert_2^2>0$,
and the strictly positive one-excursion kernel makes $\Phi_m>0$.
Combining this with
\[
 t(C_m,W)-\bigl(p^m-pq^{m-1}\bigr)\ge \Phi_m
\]
proves all the assertions.
\end{proof}

The preceding calculation also identifies the elementary certificates of
Section~2 as low-order Schur-complement expansions.  For example, since
\[
 T_Ug=\lVert g\rVert_2^2\mathbf 1+Ag,
\]
the formula for $\Phi_5$ above can be rearranged as
\[
 \Phi_5
 =4\left\lVert T_Ug+\frac{8q-5}{8}g\right\rVert_2^2
  +\left(8q^2-10q+\frac{55}{16}\right)\lVert g\rVert_2^2,
\]
which is exactly the completion of a square used in Section~2.3.
Similarly, if
\[
 s_j=\int \lambda^j\,d\mu(\lambda),
 \qquad s_0=\mu(\mathbb R),
\]
then integrating the two- and three-excursion kernels for $m=7$ gives
\[
 \begin{aligned}
 &\iint P_{7,2}(q;\lambda_1,\lambda_2)
 \,d\mu(\lambda_1)d\mu(\lambda_2)
 +6s_0^3\\
 &\qquad
 =6s_1^2+\int\bigl(s_0B_q(\lambda)+6s_0^2\bigr)\,d\mu(\lambda),
 \end{aligned}
\]
where
\[
 B_q(\lambda)
 =12\lambda^2+(36q-21)\lambda+36q^2-42q+14.
\]
Together with $P_{7,1}=P_q$, this is precisely the representation used in
Section~2.4.

Thus the short-cycle proofs are not separate phenomena: they are the first
three low-degree instances of the same Schur-complement correction.  What is
special for $m=3,5,7$ is that the resulting low-degree kernels admit immediate
sum-of-squares certificates.  For general odd $m$, direct multivariable
completion of squares is no longer practical.  The next subsection replaces
it by a uniform positive diagonalization through a Dirichlet average.
\end{document}